% GENERATED FILE. Edit canonical lessons, then run npm run generate:books.
% canonical-source-hash: fnv1a64:f11ee796a527e57d
% canonical-lessons: GE-C116-wecker, GE-C116-handy, GE-C116-kamera, GE-C116-brille, GE-C116-regenschirm

\chapter{Alarm clock, Mobile phone, Camera, Glasses, Umbrella}
\label{ch:ge-alarm-clock-mobile-phone-camera-glasses-umbrella}
\hlchaptermodality{116}

\begin{chapteropening}
\textbf{By the end of this chapter:} \emph{I can say an alarm clock, a mobile phone, a camera, glasses and an umbrella.}

Complete the last lesson of chapter 116: I can say an alarm clock, a mobile phone, a camera, glasses and an umbrella.
\end{chapteropening}

\section[der Wecker]{der Wecker — an alarm clock}

\label{lesson:GE-C116-wecker}



\begin{quote}
Before the new one: say the German for a washing machine, then the German for an armchair.

An alarm clock is for the morning: say \emph{der Morgen}, then \emph{der Tag}.
\end{quote}

\subsection*{You'll want to know: der Wecker}
\textbf{der Wecker} — \textquotedblleft{}an alarm clock\textquotedblright{}.

\begin{cousinweb}
\textbf{Wecker} is the verb \textquotedblleft{}to wake someone\textquotedblright{} + \emph{-er}: \textquotedblleft{}the waker\textquotedblright{}. The verb $\leftarrow$ Old High German \emph{wecchen} $\leftarrow$ Proto-Germanic \emph{*wakjaną}, a cousin of English \emph{wake}. (EtymDB 2.1)
\end{cousinweb}

\begin{grammarlens}[title={What you've built}]
One more word for this chapter.
\end{grammarlens}

\subsection*{Your turn}
\begin{itemize}
\raggedright
  \item \emph{Say it:} \emph{der Wecker}
  \item \emph{Say it:} \emph{der Wecker}, once more
  \item \emph{Say it:} say \emph{der Wecker}
  \item \emph{Recall:} say the German for a washing machine, then the German for an armchair, then say \emph{der Wecker} again
\end{itemize}

\begin{tcolorbox}[breakable,colback=teal!4,colframe=teal!35!black,title={Before you move on}]
What does der Wecker mean? (\textquotedblleft{}An alarm clock\textquotedblright{}.) Say it once more.
\end{tcolorbox}
\section[das Handy]{das Handy — a mobile phone}

\label{lesson:GE-C116-handy}



\begin{quote}
Before the new one: say the German for an armchair, then the German for an alarm clock.
\end{quote}

\subsection*{You'll want to know: das Handy}
\textbf{das Handy} — \textquotedblleft{}a mobile phone\textquotedblright{}.

\begin{cousinweb}
\textbf{Handy} is taken from English \emph{handy} (EtymDB 2.1), but German made it a noun for the mobile phone, a use English does not have.
\end{cousinweb}

\begin{grammarlens}[title={What you've built}]
One more word for this chapter.
\end{grammarlens}

\subsection*{Your turn}
\begin{itemize}
\raggedright
  \item \emph{Say it:} \emph{das Handy}
  \item \emph{Say it:} \emph{das Handy}, once more
  \item \emph{Say it:} say \emph{das Handy}
  \item \emph{Recall:} say the German for an armchair, then the German for an alarm clock, then say \emph{das Handy} again
\end{itemize}

\begin{tcolorbox}[breakable,colback=teal!4,colframe=teal!35!black,title={Before you move on}]
What does das Handy mean? (\textquotedblleft{}A mobile phone\textquotedblright{}.) Say it once more.
\end{tcolorbox}
\section[die Kamera]{die Kamera — a camera}

\label{lesson:GE-C116-kamera}



\begin{quote}
Before the new one: say the German for an alarm clock, then the German for a mobile phone.
\end{quote}

\subsection*{You'll want to know: die Kamera}
\textbf{die Kamera} — \textquotedblleft{}a camera\textquotedblright{}.

\begin{cousinweb}
\textbf{Kamera} $\leftarrow$ New Latin \emph{camera obscura}, \textquotedblleft{}dark chamber\textquotedblright{} $\leftarrow$ Latin \emph{camera}, \textquotedblleft{}vaulted room\textquotedblright{} $\leftarrow$ Greek \emph{kamára} (EtymDB 2.1) — borrowed: the first cameras were dark rooms.
\end{cousinweb}

\begin{grammarlens}[title={What you've built}]
One more word for this chapter.
\end{grammarlens}

\subsection*{Your turn}
\begin{itemize}
\raggedright
  \item \emph{Say it:} \emph{die Kamera}
  \item \emph{Say it:} \emph{die Kamera}, once more
  \item \emph{Say it:} say \emph{die Kamera}
  \item \emph{Recall:} say the German for an alarm clock, then the German for a mobile phone, then say \emph{die Kamera} again
\end{itemize}

\begin{tcolorbox}[breakable,colback=teal!4,colframe=teal!35!black,title={Before you move on}]
What does die Kamera mean? (\textquotedblleft{}A camera\textquotedblright{}.) Say it once more.
\end{tcolorbox}
\section[die Brille]{die Brille — glasses}

\label{lesson:GE-C116-brille}



\begin{quote}
Before the new one: say the German for a mobile phone, then the German for a camera.
\end{quote}

\subsection*{You'll want to know: die Brille}
\textbf{die Brille} — \textquotedblleft{}glasses\textquotedblright{}.

\begin{cousinweb}
\textbf{Brille} $\leftarrow$ Middle High German \emph{brille} (EtymDB 2.1), from \emph{berille}, \textquotedblleft{}beryl\textquotedblright{}, $\leftarrow$ Latin \emph{bēryllus}: the first eyeglass lenses were ground from beryl. Borrowed.
\end{cousinweb}

\begin{grammarlens}[title={What you've built}]
One more word for this chapter.
\end{grammarlens}

\subsection*{Your turn}
\begin{itemize}
\raggedright
  \item \emph{Say it:} \emph{die Brille}
  \item \emph{Say it:} \emph{die Brille}, once more
  \item \emph{Say it:} say \emph{die Brille}
  \item \emph{Recall:} say the German for a mobile phone, then the German for a camera, then say \emph{die Brille} again
\end{itemize}

\begin{tcolorbox}[breakable,colback=teal!4,colframe=teal!35!black,title={Before you move on}]
What does die Brille mean? (\textquotedblleft{}Glasses\textquotedblright{}.) Say it once more.
\end{tcolorbox}
\section[der Regenschirm]{der Regenschirm — an umbrella}

\label{lesson:GE-C116-regenschirm}



\begin{quote}
Before the new one: say the German for a camera, then the German for glasses.
\end{quote}

\subsection*{You'll want to know: der Regenschirm}
\textbf{der Regenschirm} — \textquotedblleft{}an umbrella\textquotedblright{}.

\begin{cousinweb}
\textbf{Regenschirm} is \emph{Regen}, \textquotedblleft{}rain\textquotedblright{}, + \emph{Schirm}, \textquotedblleft{}screen, shield\textquotedblright{} (EtymDB 2.1): a shield against rain.
\end{cousinweb}

\begin{grammarlens}[title={What you've built}]
One more word for this chapter.
\end{grammarlens}

\subsection*{Your turn}
\begin{itemize}
\raggedright
  \item \emph{Say it:} \emph{der Regenschirm}
  \item \emph{Say it:} \emph{der Regenschirm}, once more
  \item \emph{Say it:} say \emph{der Regenschirm}
  \item \emph{Recall:} say the German for a camera, then the German for glasses, then say \emph{der Regenschirm} again
\end{itemize}

\begin{tcolorbox}[breakable,colback=teal!4,colframe=teal!35!black,title={Before you move on}]
What does der Regenschirm mean? (\textquotedblleft{}An umbrella\textquotedblright{}.) Say it once more.
\end{tcolorbox}
